\section{Boundaries and Limitations}
\label{sec:boundaries}

We state the scope plainly.

\begin{enumerate}
\item \textbf{Testbed scope.} All claims are bounded to bearing vibration from one pair of operating conditions of one dataset (37.5\,Hz/11\,kN and 40\,Hz/10\,kN; speed and load both change), binary classification, GPU execution (no edge/real-time claims). Within each direction the training and evaluation bearings are disjoint; the reverse direction uses the same eight bearings with their roles swapped, has a single IR bearing in its evaluation set (the forward direction a single IR bearing in its training set) and so tests direction rather than new bearings; the third condition of the dataset has no IR bearing. The seeds do not resample bearings (Section~\ref{sec:protocol:testbed}), so the intervals carry no bearing-level uncertainty. Further independent datasets and multi-class settings are future work.
\item \textbf{Attribution level.} Credit assigned to ``the rest of the composition'' is by elimination within the enumerated component set; pairwise interaction terms are not decomposed.
\item \textbf{Transplant results are earlier-protocol and budget-conditional.} The graft onto S4D was run only under the earlier noise protocol, with a matched 50-epoch budget; a longer or graft-tuned budget might recover more, and we make no claim beyond \emph{this gate branch, in this host, under a matched training budget} (\ref{sec:results:graft}).
\item \textbf{Matched-SNR scope.} Training and evaluation share the same injected SNR throughout (Section~\ref{sec:protocol:testbed}); ``robustness'' in this paper is matched-noise robustness under noise-aware training, not clean-train/noisy-test generalisation. The gate's benefit may operate by exploiting matched-noise training (consistent with \ref{sec:results:perclass}) rather than by test-time denoising; attribution under mismatched-noise protocols is untested and may differ.
\item \textbf{Coverage by protocol.} Table~\ref{tab:coverage} states which operations were run under independent noise. Not run under independent noise: the graft onto S4D, the stem and $A$ ablations, additive substitution on the reverse transition, the per-class grid and the disjoint-bearing Paderborn setting. Not run at all: gate removal and additive substitution in the fully frozen block on the reverse transition, and the fully frozen block, additive arm and graft on Paderborn shared bearings.
\item \textbf{Noise-regime coverage.} Under strong coloured noise at $-10$\,dB a TCN baseline slightly overtakes the Mamba-style block, and $-10$\,dB is near-chance for the time-invariant baseline and excluded from adjudication; the demonstrated advantage regime is matched, Gaussian-dominated noise.
\item \textbf{Conv-stem dual reading.} Its clean-data cost admits both representation-mistransfer and optimisation-difficulty readings, reported in place (\ref{sec:results:localise}).
\item \textbf{Reporting rule.} Reported values are final-epoch macro-F1 under a fixed 50-epoch budget. The best-epoch rule of the original grids, for which the early thresholds were pre-specified, is optimistic by $0.6$--$10.6\pp$ (\ref{sec:results:sensitivity}). The final-epoch re-run of the original-design arms (earlier noise protocol) and the independent-noise re-run cover different sets of cells (Table~\ref{tab:coverage}); levels outside the adjudication band and the per-class grid were not re-run. No source-domain validation rule was available, and because the final epoch of the same evaluation condition has been inspected repeatedly, final-epoch values are not a fresh confirmatory test.
\item \textbf{Scope of the freezing intervention.} The first freeze arm removes the input-dependence of $\Delta$, $A$, $B$, $C$ only; the fully input-independent scan is a separate grid that was specified after the earlier grids had been seen. No equivalence margin was pre-specified, so ``no measurable cost'' means the point estimates favour the frozen arms with intervals including zero, not equivalence. The value stream and the gate input remain input-dependent, and freezing is implemented inside this Mamba-3 layer; a conclusion about input-dependence in general would need other blocks and tasks.
\item \textbf{Gate interventions.} Under independent noise, gate removal is resolved only at $-6$\,dB in the partially and the fully frozen block on the original transition, and without its gate the fully frozen block still exceeds S4D at $0$ and $-2$\,dB; the combined statement that the fully frozen block's advantage is carried by the gate is therefore not made. Where the gate-removal cost is resolved at every level (partially frozen block, reverse transition) it was not tested in the fully frozen block. The additive substitution could not be distinguished from gate removal, but the multiplicative--additive differences are positive at $-6$\,dB; with gate removal only weakened the pre-specified rescue rule is not assessable, and neither equivalence nor zero rescue is claimed. The rescue ratio is a conditional summary (levels with a denominator below $2\pp$ are excluded in some resamples).
\item \textbf{Noise protocol.} All grids not marked in Table~\ref{tab:coverage} as independent-noise used training and evaluation noise shared by index. Where both protocols exist the arm means differ by at most $4\pp$ and no contrast changed direction, which is not a guarantee for the arms that were not re-run.
\item \textbf{Strength of the pre-specification evidence.} Thresholds were pre-specified internally; for the follow-up grids from the width-matched graft onwards the thresholds were also committed to a public repository before the first cell, for the earlier grids the evidence is local file times and run-directory names, not an external time-stamped registry (Section~\ref{sec:protocol:prereg}). The noise-depth analysis (Section~\ref{sec:results:depth}) and the multiplicative--additive differences were added during review and are not pre-specified. The process-level repeat check is a single observation.
\item \textbf{Construction checks.} For the frozen arm the state-dictionary check compares names and shapes of non-target tensors; value-level comparison was applied to the component-ablation and graft arms only.
\item \textbf{Second dataset.} On Paderborn, the gate-removal cost replicates when training and evaluation share bearings under independent noise, and the partially frozen block's advantage over S4D is resolved at $-6$\,dB only. With bearings disjoint from the training set (earlier protocol only; per-index label agreement 0.999) nothing replicates and every architecture, S4D included, transferred poorly (47.7--56.9\% against a balanced-chance level of about 50\%). The two Paderborn settings differ in more than bearing overlap, so they do not identify why transfer fails. The fully frozen block was run only in the disjoint setting.
\item \textbf{Sample size.} Every cell mean uses five seeds. The smallest attainable one-sided Wilcoxon $p$ is $0.031$, paired intervals are wide, and several secondary statements (conv stem, $A$ parameterisation, gateless versus S4D) are not resolved at this size. For the original-design grids conclusions rest on effect sizes that clear the pre-specified thresholds; for the later grids they rest on paired intervals as defined in Section~\ref{sec:protocol:prereg}. Seeds resample initialisation, batch order and noise (evaluation noise is one bank per seed), not bearings.
\end{enumerate}

A single-repeat check of process-level non-determinism ($\sim\!0.7\pp$) and the label-hash/adjudication trail are documented in the reproducibility appendices (\ref{app:fairness}--\ref{app:confusion}).
